\section{Auditing and Mitigation}

\begin{frame}{The Audit Loop}

    \centering
    \resizebox{\textwidth}{!}{%
    \begin{tikzpicture}[
        node distance=0.65cm,
        auditstep/.style={rectangle, draw, rounded corners, minimum height=1.15cm, text width=2.35cm,
                     align=center, font=\scriptsize\bfseries, fill=raiBlue!10},
        arr/.style={-{Latex[length=2.2mm]}, thick}
    ]
        \node[auditstep] (a) {1. Name the harm and the affected groups};
        \node[auditstep, right=of a] (b) {2. Choose the metric(s) and the criterion};
        \node[auditstep, right=of b] (c) {3. Measure per slice, with sample sizes};
        \node[auditstep, right=of c] (d) {4. Diagnose \emph{why} the gap exists};
        \node[auditstep, right=of d] (e) {5. Mitigate and re-measure};
        \node[auditstep, right=of e, fill=raiGreen!12] (f) {6. Document and monitor};

        \foreach \x/\y in {a/b, b/c, c/d, d/e, e/f} \draw[arr] (\x) -- (\y);
        \draw[arr, dashed] (f.south) .. controls +(0,-1.0) and +(0,-1.0) .. (c.south);
    \end{tikzpicture}}

    \vspace{1.2em}

    \raggedright
    \begin{itemize}
        \setlength\itemsep{0.45em}
        \item Step 4 is the one people skip. A gap caused by \textbf{too few samples} needs data;
              a gap caused by a \textbf{proxy label} needs a new label; a gap caused by a
              \textbf{threshold} needs a threshold. The fixes are not interchangeable.
        \item Step 6 is not optional: distributions move, and a model that was audited in March is
              not audited in December.
    \end{itemize}
\end{frame}

\begin{frame}{Disaggregated Evaluation: The Core Technique}

    \textit{One number for the test set, then the same number for every slice you can define.}

    \vspace{0.7em}
    \centering
    \scriptsize
    \begin{tabular}{lrrrrr}
        \toprule
        \textbf{Slice} & \textbf{$n$} & \textbf{Selection rate} & \textbf{TPR} & \textbf{FPR} & \textbf{PPV} \\
        \midrule
        Overall            & 10{,}000 & 0.24 & 0.78 & 0.11 & 0.71 \\
        \midrule
        Group A            &  8{,}200 & 0.26 & 0.81 & 0.10 & 0.73 \\
        Group B            &  1{,}800 & 0.14 & 0.62 & 0.13 & 0.58 \\
        \midrule
        Group B, age $<30$ &    420   & 0.11 & 0.49 & 0.14 & 0.51 \\
        Group B, age $\geq 60$ &  180 & 0.09 & 0.44 & 0.12 & 0.47 \\
        \bottomrule
    \end{tabular}

    \vspace{0.5em}
    {\scriptsize Illustrative numbers, constructed for this slide.}

    \vspace{0.8em}
    \raggedright
    \begin{itemize}\scriptsize
        \setlength\itemsep{0.35em}
        \item The overall row is fine. Group B is not, and the intersectional rows are worse ---
              the pattern from Gender Shades, reproduced in miniature.
        \item Always print \textbf{$n$} beside the metric. Always print the \textbf{worst slice}
              in the summary, not just the mean.
        \item In \texttt{scikit-learn}, this is a \texttt{groupby} plus
              \texttt{classification\_report}; in \texttt{fairlearn} it is \texttt{MetricFrame}.
    \end{itemize}
\end{frame}

\begin{frame}{Four Things That Go Wrong in Practice}

    \begin{itemize}
        \setlength\itemsep{0.65em}
        \item \textbf{Small slices are noisy.} A TPR of $0.44$ estimated from 180 positive cases
              carries a 95\% interval of about $\pm 0.07$; from 40 cases, about $\pm 0.15$. Report
              intervals, or you will chase noise and ``fix'' gaps that were never there.
              A bootstrap is fine here.

        \item \textbf{Multiple comparisons.} Slice by 8 attributes and some slice will look bad by
              chance. Pre-register the slices you care about; treat the rest as exploratory.

        \item \textbf{You may not have the group labels.} They are often unavailable by law or by
              policy. Options: a consented subset for audit only, an
              independently collected audit set, or proxy-based estimates --- each with its own
              error, which must be stated.

        \item \textbf{Intersections multiply.} Four binary attributes give 16 cells and most will
              be tiny. Choose the intersections that correspond to a plausible harm, rather than
              enumerating everything.
    \end{itemize}
\end{frame}

\begin{frame}{Where Mitigation Can Be Applied}

    \centering
    \begin{tikzpicture}[
        node distance=1.0cm,
        stage/.style={rectangle, draw, rounded corners, minimum height=1.0cm, text width=2.4cm,
                      align=center, font=\small\bfseries},
        note/.style={align=center, font=\scriptsize, text width=3.1cm},
        arr/.style={-{Latex[length=2.5mm]}, thick}
    ]
        \node[stage, fill=raiBlue!12]  (data)  {Data};
        \node[stage, fill=raiAmber!15, right=of data] (train) {Training};
        \node[stage, fill=raiGreen!12, right=of train] (out)  {Outputs};

        \draw[arr] (data) -- (train);
        \draw[arr] (train) -- (out);

        \node[note, below=0.45cm of data]  {\textbf{Pre-processing}\\ reweigh, resample, transform features};
        \node[note, below=0.45cm of train] {\textbf{In-processing}\\ fairness constraint or penalty in the objective};
        \node[note, below=0.45cm of out]   {\textbf{Post-processing}\\ adjust thresholds or scores per group};
    \end{tikzpicture}

    \vspace{1.4em}

    \raggedright
    \begin{itemize}\small
        \setlength\itemsep{0.35em}
        \item \textbf{Access decides the option.} Third-party model behind an API: post-processing
              only. Your own training loop: all three.
        \item \textbf{None of them repair a bad label or a missing population.} Those are data
              problems, and mitigation applied on top will only redistribute the error.
    \end{itemize}
\end{frame}

\begin{frame}{Pre-processing and In-processing}

    \begin{columns}[T]
        \begin{column}{0.49\textwidth}
            \textbf{Pre-processing} --- change the data so the disparity is not there to be learned.
            \vspace{0.4em}
            \begin{itemize}\scriptsize
                \setlength\itemsep{0.32em}
                \item \textbf{Reweighing} (Kamiran \& Calders, 2012): weight each
                      $(\text{group},\text{label})$ cell so group and label become independent in
                      the weighted sample. Works with any estimator taking
                      \texttt{sample\_weight}.
                \item \textbf{Resampling}: same effect, but duplicates or discards rows.
                \item \textbf{Targeted data collection}: usually the highest-value intervention,
                      and the one most often skipped because it costs money rather than code.
                \item \textbf{Representation transforms}: learn an encoding from which $A$ cannot
                      be predicted.
            \end{itemize}
            \vspace{0.4em}
            {\scriptsize Model-agnostic and done once --- but you optimise a proxy, and the
            downstream model can partly undo it.}
        \end{column}
        \begin{column}{0.49\textwidth}
            \textbf{In-processing} --- put the requirement into the training objective.
            \vspace{0.4em}
            \begin{itemize}\scriptsize
                \setlength\itemsep{0.32em}
                \item \textbf{Constrained optimisation}: minimise loss subject to
                      $|\text{TPR}_\GrpA - \text{TPR}_\GrpB| \leq \varepsilon$. Agarwal et al.\ (2018)
                      reduce this to a sequence of ordinary cost-sensitive fits, so any learner
                      works --- this is \texttt{fairlearn}'s \texttt{ExponentiatedGradient}.
                \item \textbf{Regularisation}: add $\lambda \cdot (\text{disparity})$ to the loss;
                      $\lambda$ is an explicit dial on the trade-off.
                \item \textbf{Adversarial debiasing}: train the predictor against an adversary
                      trying to recover $A$ from the prediction or representation.
            \end{itemize}
            \vspace{0.4em}
            {\scriptsize Enforced where the model is fitted --- but needs training access, adds
            tuning, and holds on the \emph{training} distribution.}
        \end{column}
    \end{columns}

    \vspace{0.4em}
    {\scriptsize F.~Kamiran and T.~Calders, \emph{Knowledge and Information Systems} 33(1), 2012.
    A.~Agarwal, A.~Beygelzimer, M.~Dud\accent19 ik, J.~Langford and H.~Wallach, ``A Reductions
    Approach to Fair Classification,'' ICML 2018,
    \href{https://arxiv.org/abs/1803.02453v3}{arXiv:1803.02453v3}.}
\end{frame}

\begin{frame}{Post-processing: Group-Specific Thresholds}

    \begin{columns}[T]
        \begin{column}{0.52\textwidth}
            Hardt et al.\ (2016) show that if you keep the score $R$ and choose a
            \textbf{separate threshold per group}, you can achieve equalised odds --- and they give
            the optimal choice.

            \vspace{0.6em}
            \begin{itemize}\scriptsize
                \setlength\itemsep{0.3em}
                \item Each group's ROC curve gives the achievable (FPR, TPR) pairs.
                \item Equalised odds requires landing on the \textbf{same point} for both groups,
                      so the feasible set is the intersection.
                \item The best you can do is the intersection point that maximises utility ---
                      which is generally \emph{worse} than either group's own optimum.
            \end{itemize}
        \end{column}
        \begin{column}{0.46\textwidth}
            \centering
            \begin{tikzpicture}[scale=2.6]
                \draw[->] (0,0) -- (1.12,0) node[right, font=\scriptsize] {FPR};
                \draw[->] (0,0) -- (0,1.12) node[above, font=\scriptsize] {TPR};
                \draw[gray!50, dashed] (0,0) -- (1,1);
                \draw[raiBlue, very thick] (0,0) .. controls (0.08,0.72) and (0.42,0.98) .. (1,1)
                    node[pos=0.55, above left, font=\scriptsize, text=raiBlue] {Group A};
                \draw[raiRed, very thick]  (0,0) .. controls (0.28,0.46) and (0.66,0.86) .. (1,1)
                    node[pos=0.62, below right, font=\scriptsize, text=raiRed] {Group B};
                \fill[black] (0.478,0.62) circle (0.018)
                    node[right=2pt, font=\scriptsize] {reachable by both};
            \end{tikzpicture}
            \vspace{0.2em}
            {\scriptsize Schematic drawn for this slide, after the construction in Hardt et al.\ (2016).}
        \end{column}
    \end{columns}

    \vspace{0.6em}
    \begin{tcolorbox}[colback=raiAmber!5, colframe=raiAmber!70, title=Legal caveat]
        \scriptsize Applying a different threshold by protected group can itself be unlawful in some
        jurisdictions and domains --- for example, adjusting employment-test scores by race is
        prohibited in the United States under the Civil Rights Act of 1991. Check before you ship;
        the statistically optimal fix is not always the permitted one.
    \end{tcolorbox}
\end{frame}

\begin{frame}{Trade-offs, Monitoring, and Tooling}

    \begin{columns}[T]
        \begin{column}{0.5\textwidth}
            \textbf{Expect to give something up}
            \begin{itemize}\small
                \setlength\itemsep{0.3em}
                \item Constraining a metric usually costs some accuracy. Report the curve
                      (disparity vs accuracy), not one point --- it turns an argument into a
                      decision.
                \item Sometimes the cost is near zero, which is worth discovering.
                \item ``Levelling down'' --- degrading the strong group to close the gap ---
                      satisfies the metric and helps nobody. Say explicitly whether it is acceptable.
            \end{itemize}
        \end{column}
        \begin{column}{0.48\textwidth}
            \textbf{After deployment}
            \begin{itemize}\small
                \setlength\itemsep{0.3em}
                \item Re-run the disaggregated report on production traffic on a schedule.
                \item Alert on disparity, not only on accuracy drift.
                \item Keep a channel for people to contest a decision, and log contested cases ---
                      they are your best source of failure modes.
            \end{itemize}
            \vspace{0.5em}
            \textbf{Libraries}
            \begin{itemize}\small
                \item \texttt{fairlearn} (metrics, reductions, threshold optimiser)
                \item \texttt{AIF360} (a wide catalogue of algorithms)
                \item \texttt{aequitas} (audit reports)
            \end{itemize}
        \end{column}
    \end{columns}

    \vspace{0.7em}
    \begin{block}{}
        \centering
        \small \textbf{Exercise} (companion notebook in preparation): disaggregate on UCI Adult,
        show the parity/odds tension, apply threshold post-processing, and measure what it cost.
    \end{block}
\end{frame}
